% Three ways to run the same hierarchy: only the endpoints change.
\begin{tikzpicture}[font=\sffamily\scriptsize, >=Latex,
  col/.style={draw=octrule!40, rounded corners=3pt, inner sep=0, minimum width=46mm, minimum height=64mm, anchor=north},
  src/.style={draw=octamber, fill=octctl, rounded corners=2pt, text width=38mm, minimum height=11mm, align=center, inner sep=2pt},
  mid/.style={draw=octrule, fill=octmem, rounded corners=2pt, text width=38mm, minimum height=22mm, align=center, inner sep=2pt},
  snk/.style={draw=octrule, fill=white, rounded corners=2pt, text width=38mm, minimum height=11mm, align=center, inner sep=2pt},
  ttl/.style={font=\sffamily\footnotesize\bfseries, text=octink},
  w/.style={<->, thick, octink}]
  \foreach \ox/\ot/\os/\osb/\om/\omb in {
      0/{Standalone}/{\textbf{CPU}: trace replay}/{trace\_C<n>.trc.shared, or a periodic task}/{\textbf{MainMemoryController}}/{fixed 250-cycle latency},
      50/{Standalone + MCsim}/{\textbf{CPU}: trace replay}/{the same traces}/{\textbf{MCsimInterface}}/{DDR4, twenty schedulers},
      100/{Full system}/{\textbf{ExternalCPU}: gem5}/{Arm cores, Linux, real applications; ATP aggressors}/{\textbf{MCsimInterface}}/{DDR4; gem5 holds the data}} {
    \node[col] (c) at (\ox mm, 0) {};
    \node[ttl] at ($(c.north)+(0,-4mm)$) {\ot};
    \node[src] (s) at ($(c.north)+(0,-14mm)$) {\os\\\scriptsize\osb};
    \node[mid] (m) at ($(c.north)+(0,-34mm)$) {\textbf{the same hierarchy}\\L1s \textperiodcentered\ bus[0] \textperiodcentered\ LLC \textperiodcentered\ bus[1]\\\scriptsize CSV-configured, unchanged};
    \node[snk] (k) at ($(c.north)+(0,-54mm)$) {\om\\\scriptsize\omb};
    \draw[w] (s) -- (m); \draw[w] (m) -- (k);
  }
  \node[text=octquiet, align=center, text width=130mm] at (50mm, -70mm) {the endpoints are \code{CommunicationInterface} producers and consumers; swapping them changes the request source and the memory sink, nothing in between};
\end{tikzpicture}
